% Lösungen Einheit 3 von 3 – Unabhängigkeit als Argument: die Ausgleichs-Fehlvorstellung widerlegen (zusammenbau v0.1)
\einheitenkopf{Einheit 3 von 3 · Unabhängigkeit als Argument: die Ausgleichs-Fehlvorstellung widerlegen}

\erg{14}{a) $0{,}5$ – die Münze hat kein Gedächtnis, jeder Wurf ist unabhängig \quad b) Falsch: Die Wahrscheinlichkeit für Blau ist bei jeder Drehung gleich groß und hängt nicht von früheren Drehungen ab; auch bei den nächsten Drehungen ist etwa ein Viertel Blau zu erwarten, ein Ausgleich ist nicht zu erwarten. \quad c) Falsch: Die Samstage sind voneinander unabhängig, die Wahrscheinlichkeit für sieben richtige Antworten ist an jedem Samstag dieselbe; die früheren Ergebnisse ändern sie nicht – dass vier solche Samstage nacheinander selten sind, ist kein Argument für den einzelnen Samstag.}
\erg{15}{a) Fehler: Das Gesetz der großen Zahlen verspricht keinen Ausgleich, sondern nur, dass die relative Häufigkeit bei sehr vielen Würfen nahe der Wahrscheinlichkeit liegt – durch das Gewicht der vielen weiteren Würfe. Richtig: Jeder weitere Wurf zeigt Wappen mit der Wahrscheinlichkeit ein halb, unabhängig von den bisherigen. \quad b) Die Würfe sind unabhängig: Die Münze verändert sich durch einen Wurf nicht, und das Ergebnis eines Wurfs wirkt nicht auf den nächsten – die Wahrscheinlichkeit ist bei jedem Wurf dieselbe.}
